\begin{textsample}{POS Dim 1 – 2021\_c5 – Score 18.00 – 2021\_c5/enade\_2021\_geo\_b7744036e1.txt}  \label{ex:f1_pos_038}
CAPÍTULO 5 DISTRIBUIÇÃO DOS CONCEITOS 5.1 \textbf{PANORAMA} NACIONAL DA DISTRIBUIÇÃO DOS CONCEITOS O Conceito Enade é um indicador de qualidade que avalia os cursos por \textbf{intermédio} dos desempenhos dos estudantes a partir dos resultados do Enade.
Os valores possíveis vão de 1 ( pior situação ) a 5 ( melhor situação ).
Os cursos com apenas um ou sem nenhum concluinte participante não obtêm o Conceito Enade, ficando Sem Conceito ( SC ).
Foram avaliados em Geografia 299 cursos, sendo 229 ( 76,6 \% ) de Licenciatura e 70 ( 23,4 \% ) de Bacharelado.
Enfatiza -se, mais uma vez, que a diferença entre os cursos \textbf{tabulados} neste capítulo e no Capítulo 2 corresponde aos cursos que não foram avaliados, ou seja, àqueles, em princípio, sem alunos concluintes que tenham sido inscritos para a prova.
Na Tabela 5.1a, são apresentadas a quantidade e a distribuição de cursos de Licenciatura de Geografia participantes do Enade/2021, por faixa de conceito e Grande Região.
Observando -se os dados da Tabela 5.1a, nota -se que, dos 229 cursos participantes, 82 ( 35,8 \% ) classificaram -se com conceito 3, o valor modal.
Esse foi também o conceito modal nas regiões Nordeste ( 42,7 \% ) e Sul ( 51,4 \% ).
O conceito 2 foi o segundo mais frequente em nível nacional ( 31,4 \%, correspondendo a 72 cursos ), tendo sido o conceito modal das regiões Norte ( 50,0 \% ) e Centro-Oeste ( 45,2 \% ), e o conceito 4, o terceiro ( 21,8 \%, correspondendo a 50 cursos ), tendo sido o conceito modal na região Sudeste ( 41,2 \% ).
Houve, ainda, nove cursos ( 3,9 \% ) que receberam o conceito 1, e 11 cursos ( 4,8 \% ) que receberam o conceito 5.
Dos 229 cursos de Geografia, cinco ( 2,2 \% ) ficaram Sem Conceito ( SC ).
A região Norte participou com 18 cursos, o correspondente a 7,9 \% do total nacional.
Desses, nove cursos receberam o conceito 2, ( o conceito modal na região, como já comentado ), o que \textbf{equivale} a 50,0 \% do total regional.
A três cursos ( 16,7 \% ) atribuiu -se o conceito 1, a quatro cursos ( 22,2 \% ), o conceito 3, e a um curso ( 5,6 \% ), o conceito 5.
Um curso ficou Sem Conceito ( SC ), e nenhum curso recebeu o conceito 5 nessa região.
A região Nordeste participou com 75 cursos, o equivalente a 32,8 \% do total nacional.
Nessa região, 32 cursos ( 42,7 \% em termos regionais ) obtiveram conceito 3, o conceito modal para a região, como já comentado.
Os conceitos 1 e 2 foram atribuídos, respectivamente, a três e a 28 cursos ( 4,0 \% e 37,3 \%, respectivamente ).
Ainda nessa região, dez cursos ( 13,3 \% ) foram avaliados com o conceito 4, e dois cursos ( 2,7 \% ) foram avaliados com o conceito 5.
Nenhum curso ficou Sem Conceito ( SC ).
Dos 68 cursos participantes da região Sudeste ( 29,7 \% do total nacional ), o conceito 4 foi o conceito modal, como já anteriormente assinalado, tendo sido obtido por 28 cursos ( 41,2 \% ).
O conceito 1 foi atribuído a um curso ( 1,5 \% ), e o conceito 2, a 14 cursos ( 20,6 \% ).
Receberam os conceitos 3 e 5, respectivamente, 17 cursos ( 25,0 \% ) e cinco cursos ( 7,4 \% ).
Nessa região, três cursos ficaram Sem Conceito ( SC ).
Os 37 cursos da região Sul corresponderam a 16,2 \% do total nacional.
Como já \textbf{indicado}, houve predominância do conceito 3, atribuído a 51,4 \% dos cursos da região ( conceito modal ), ou seja, atribuído a 19 dos 37 cursos participantes na região Sul.
O conceito 2 foi atribuído a sete cursos ( 18,9 \% ), e os conceitos 4 e 5, a sete ( 18,9 \% ) e a três cursos ( 8,1 \% ), respectivamente.
Um curso ficou Sem Conceito ( SC ) na região Sul, e nenhum recebeu conceito 1.
Já dos 31 cursos participantes na região Centro-Oeste ( 13,5 \% do total nacional ), 14 ( 45,2 \% em termos regionais ) receberam conceito 2, o conceito modal, como já destacado.
Dois cursos ( 6,5 \% ) obtiveram o conceito 1, dez cursos ( 32,3 \% ) obtiveram o conceito 3, e outros quatro ( 12,9 \% ), o conceito 4.
Apenas um curso obteve o conceito 5.
Nessa região, nenhum curso ficou Sem Conceito ( SC ).
No Gráfico 5.1a, é apresentada a distribuição \textbf{cumulativa} do Conceito Enade dos cursos de Licenciatura, segundo a Grande Região.
Em um gráfico de distribuição \textbf{cumulativa}, em cada valor de \textbf{abscissa} ( eixo x ), o valor apresentado no eixo das \textbf{ordenadas} ( eixo y ) é o valor acumulado da \textbf{variável} ( Conceito Enade ).
Observando -se a linha vertical \textbf{pontilhada}, que ajuda a \textbf{visualizar} o percentual de conceitos acumulados até o valor 3, ou seja, os conceitos 1, 2 e 3, pode -se notar que, por exemplo, a região Norte ( linha azul ) apresentou 94,1 \% dos cursos com o conceito 3 ou abaixo desse valor.
A região Sudeste ( linha vermelha ), por outro lado, apresentou menos de 50 \% dos cursos na mesma situação : com o conceito 3 ou abaixo.
Destaca -se que, em algumas regiões, o conceito 5 pode não ter sido alcançado por nenhum curso.
Foi o que ocorreu na região Norte ( linha azul ), em que todos os cursos obtiveram conceitos menores ou iguais a 4, tendo a \textbf{poligonal} alcançado 100 \% no conceito 4.
Em linhas gerais, em um gráfico de distribuição \textbf{cumulativa}, \textbf{poligonais} mais à esquerda representam as Grandes Regiões com a pior distribuição, e \textbf{poligonais} mais à direita, as Grandes Regiões com os melhores conceitos.
A região Norte ( linha azul ) apresentou os cursos com os conceitos mais baixos, ao passo que a região Sudeste ( linha vermelha ) apresentou a distribuição com os valores maiores.
Cursos em IES situadas nas regiões Nordeste ( linha verde ), Centro-Oeste ( linha laranja ) e Sul ( linha \textbf{rosa} ) apresentaram uma situação intermediária.
Na Tabela 5.1b, são apresentadas a quantidade e a distribuição de cursos de Bacharelado de Geografia participantes do Enade/2021, por faixa de conceito e Grande Região.
Observando -se os dados da Tabela 5.1b, nota -se que, dos 70 cursos participantes, 27 ( 38,6 \% ) classificaram -se com conceito 3, o valor modal.
Esse foi também o conceito modal nas regiões Norte ( 50,0 \% ), Sudeste ( 39,1 \%, sendo que o conceito 4 também foi igualmente atribuído a 39,1 \% dos cursos ), Sul ( 50,0 \% ) e Centro-Oeste ( 33,3 \%, sendo que o conceito 2 também foram atribuído a 33,3 \% dos cursos ).
O conceito 4 foi o segundo mais frequente em nível nacional ( 30,0 \%, correspondendo a 21 cursos ), e o conceito 2, o terceiro ( 18,6 \%, correspondendo a 13 cursos ).
Houve, ainda, três cursos ( 4,3 \% ) que receberam o conceito 1, e cinco cursos ( 7,1 \% ) que receberam o conceito 5.
Apenas um curso ( 1,4 \% ) ficou Sem Conceito ( SC ).
A região Norte participou com seis cursos, o correspondente a 8,6 \% do total nacional.
Desses, três cursos receberam o conceito modal 3, o que \textbf{equivale} a 50,0 \% do total regional.
A um curso ( 16,7 \% ) atribuiu -se o conceito 1, e a dois cursos ( 33,3 \% ), o conceito 2.
Nenhum curso ficou Sem Conceito ( SC ) ou recebeu os conceitos 4 ou 5.
A região Nordeste participou com 16 cursos, o equivalente a 22,9 \% do total nacional.
Nessa região, sete cursos ( 43,8 \% em termos regionais ) obtiveram conceito 2, o conceito modal para a região.
Os conceitos 3 e 4 foram atribuídos, respectivamente, a quatro e a cinco cursos ( 25,0 \% e 31,3 \%, respectivamente ).
Nenhum curso ficou Sem Conceito ( SC ) ou recebeu os conceitos 1 ou 5.
Dos 23 cursos participantes da região Sudeste ( 32,9 \% do total nacional ), os conceitos 3 e 4 foram os conceitos modais ( distribuição \textbf{bimodal} ), como já anteriormente assinalado, tendo sido obtidos por nove cursos ( 39,1 \% ) cada um.
O conceito 1 foi atribuído a um curso ( 4,3 \% ), e o conceito 2, também a um curso ( 4,3 \% ).
Receberam o conceito 5 três cursos ( 13,0 \% ).
Nessa região, nenhum curso ficou Sem Conceito ( SC ).
Os 16 cursos da região Sul corresponderam a 22,9 \% do total nacional.
Como já \textbf{indicado}, houve predominância do conceito 3, atribuído a 50,0 \% dos cursos da região ( conceito modal ), ou seja, atribuído a oito dos 16 cursos participantes na região Sul.
O conceito 4 foi atribuído a seis cursos ( 37,5 \% ), e o conceito 5, a um curso ( 6,3 \% ).
Um curso ( 6,3 \% ) ficou Sem Conceito.
Já dos nove cursos participantes na região Centro-Oeste ( 12,9 \% do total nacional ), três ( 33,3 \% em termos regionais ) receberam o conceito 2, e outros três o conceito 3, uma distribuição \textbf{bimodal}.
Os conceitos 1, 4 e 5 foram atribuídos a um curso cada um ( 11,1 \% cada ).
Nenhum curso ficou Sem Conceito ( SC ).
No Gráfico 5.1b, é apresentada a distribuição \textbf{cumulativa} do Conceito Enade dos cursos de Bacharelado, segundo a Grande Região.
Em um gráfico de distribuição \textbf{cumulativa}, em cada valor de \textbf{abscissa} ( eixo x ), o valor apresentado no eixo das \textbf{ordenadas} ( eixo y ) é o valor acumulado da \textbf{variável} ( Conceito Enade ).
Observando -se a linha vertical \textbf{pontilhada}, que ajuda a \textbf{visualizar} o percentual de conceitos acumulados até o valor 3, ou seja, os conceitos 1, 2 e 3, pode -se notar que, por exemplo, a região Norte ( linha azul ) apresentou 100 \% dos cursos com o conceito 3 ou abaixo desse valor.
A região Sudeste ( linha vermelha ), por outro lado, apresentou 47,8 \% dos cursos na mesma situação : com o conceito 3 ou abaixo.
Destaca -se que, em algumas regiões, o conceito 5 pode não ter sido alcançado por nenhum curso, como ocorreu nas regiões Norte ( linha azul ), na qual todos os cursos obtiveram conceitos menores ou iguais a 3, tendo a sua \textbf{poligonal} alcançado 100 \% no conceito 3, e Nordeste ( linha verde ), na qual todos os cursos obtiveram conceitos menores ou iguais a 4, tendo a sua \textbf{poligonal} alcançado 100 \% no conceito 4.
Em linhas gerais, em um gráfico de distribuição \textbf{cumulativa}, \textbf{poligonais} mais à esquerda representam as Grandes Regiões com a pior distribuição, e \textbf{poligonais} mais à direita, as Grandes Regiões com os melhores conceitos.
A região Norte ( linha azul ) apresentou os cursos com os conceitos mais baixos, ao passo que as regiões Sul ( linha \textbf{rosa} ) e Sudeste ( linha vermelha ) apresentaram as distribuições com os valores maiores.
Cursos em IES situadas nas regiões Nordeste ( linha verde ) e Centro-Oeste ( linha laranja ) apresentaram uma situação intermediária. 5.2 CONCEITOS POR CATEGORIA ADMINISTRATIVA, POR MODALIDADE DE ENSINO E GRANDE REGIÃO Na Tabela 5.2a, é apresentada a distribuição dos cursos dos cursos de Licenciatura de Geografia participantes do Enade/2021, por Categoria Administrativa e Modalidade de Ensino, de acordo com os conceitos por eles alcançados, segundo as Grandes Regiões brasileiras.
Dos 229 cursos participantes, 181 ( 79,0 \% ) eram \textbf{ministrados} em Instituições \textbf{Públicas}, e 48 ( 21,0 \% ), em Instituições Privadas.
Quanto à Modalidade de Ensino, 43 cursos ( 18,8 \% ) eram a Distância, e 186 ( 81,2 \% ), Presenciais.
De acordo com as informações da Tabela 5.2a, em termos nacionais, entre Instituições \textbf{Públicas} e Privadas, observa -se que, dos 11 cursos avaliados com o conceito 5, nove foram oferecidos em IES Públicas, e dois, em IES Privadas.
Dos 181 cursos participantes de IES Públicas, o conceito 3 foi o valor modal, atribuído a 63 cursos ( 34,8 \% ).
Dos demais cursos de Instituições \textbf{Públicas} participantes, oito cursos ( 4,4 \% ) foram avaliados com o conceito 1 ; 56 cursos ( 30,9 \% ), com o conceito 2 ; 43 ( 23,8 \% ), com o conceito 4 ; e nove ( 5,0 \% ), com o conceito 5.
Nessa categoria, dois cursos ficaram Sem Conceito ( SC ).
Na rede privada, o conceito modal também foi o 3, atribuído a 19 cursos, o correspondente a 39,6 \% dos 48 cursos da categoria.
Dos demais cursos participantes, um ( 2,1 \% ) recebeu o conceito 1, e 16 ( 33,3 \% ) receberam o conceito 2.
O conceito 4 foi atribuído a sete cursos ( 14,6 \% ), e o conceito 5, a dois cursos ( 4,2 \% ).
Nessa Categoria Administrativa, três cursos ( 6,3 \% ) ficaram Sem Conceito ( SC ).
Dos 186 cursos participantes da Modalidade Presencial, o conceito 3 foi o valor modal, atribuído a 66 cursos ( 35,5 \% ).
Dos demais cursos nessa Modalidade de Ensino, oito cursos ( 4,3 \% ) receberam o conceito 1 ; 52 cursos ( 28,0 \% ), o conceito 2, 46 cursos ( 24,7 \% ), o conceito 4 ; e dez cursos ( 5,4 \% ), o conceito 5.
Nessa Modalidade de Ensino, quatro cursos ( 2,2 \% ) ficaram Sem Conceito ( SC ).
Já dos 43 cursos participantes da Modalidade a Distância, 20 ( 46,5 \% ) obtiveram o conceito modal 2.
Dos demais cursos nessa Modalidade de Ensino, um curso ( 2,3 \% ) recebeu o conceito 1 ; 16 cursos ( 37,2 \% ) receberam o conceito 3 ; quatro cursos ( 9,3 \% ), o conceito 4 ; e um curso ( 2,3 \% ) obteve o conceito 5.
Um curso ( 2,3 \% ) nessa Modalidade de Ensino ficou Sem Conceito ( SC ).
Na análise por região, observa -se que, na região Norte, todos os 18 cursos participantes foram mantidos por Instituições \textbf{Públicas} e na Modalidade Presencial.
O conceito modal para esses cursos foi o 2, atribuído a nove cursos, correspondendo a 50,0 \% dessa categoria na região.
Três cursos foram avaliados com o conceito 1, quatro, com o conceito 3, e um foi avaliado com o conceito 4.
Um curso ficou Sem Conceito ( SC ) nessa combinação de categoria e região.
Na região Nordeste, dos 75 cursos participantes, dois eram da rede privada, o equivalente a 2,7 \% do total da região.
Desses, um curso foi avaliado com o conceito 2, e o outro, com o conceito 5.
As Instituições \textbf{Públicas} da região Nordeste participaram com 73 cursos ( 97,3 \% do total da região ).
Desses, 32 obtiveram conceito 3, o valor modal.
Os demais conceitos dessa categoria foram : 1 ( 3 cursos ), 2 ( 27 cursos ), 4 ( dez cursos ) e 5 ( um curso ).
Nenhum curso ficou Sem Conceito ( SC ).
Dos 65 cursos ( 86,7 \% ) participantes da Modalidade Presencial na região Nordeste, o conceito 3 foi o valor modal, atribuído a 29 cursos ( 44,6 \% ).
Dos demais cursos nessa Modalidade de Ensino, dois cursos ( 3,1 \% ) receberam o conceito 1 e 22 cursos ( 33,8 \% ), o conceito 2.
O conceito 4 foi atribuído a dez cursos ( 15,4 \% ), e o conceito 5, a dois cursos ( 3,1 \% ).
Nessa Modalidade de Ensino, nenhum curso ficou Sem Conceito ( SC ).
Já dos dez cursos ( 13,3 \% ) a Distância informados na região Nordeste, seis cursos ( 60,0 \% ) receberam o conceito modal 2, um curso ( 10,0 \% ) recebeu o conceito 1, e três cursos ( 30,0 \% ) receberam o conceito 3.
Nenhum curso nessa Modalidade de Ensino na região Nordeste ficou Sem Conceito ( SC ) ou recebeu os conceitos 4 ou 5.
Na região Sudeste, a proporção de cursos da rede privada foi de 41,2 \%, correspondendo a 28 dos 68 cursos participantes.
Nessa categoria e região, o conceito modal foi o 3, atribuído a dez cursos, e dois cursos ficaram Sem Conceito ( SC ).
Os demais cursos foram avaliados com os seguintes conceitos : 1 ( um curso ), 2 ( nove cursos ), 4 ( cinco cursos ) e 5 ( um curso ).
Dos 40 cursos ( 58,8 \% ) oferecidos por Instituições \textbf{Públicas} na região Sudeste, o conceito modal foi o 4, atribuído a 23 cursos.
Nessa categoria e nessa região, os demais cursos foram avaliados com os seguintes conceitos : 2 ( cinco cursos ), 3 ( sete cursos ) e 5 ( quatro cursos ).
Um curso ficou Sem Conceito ( SC ), e nenhum curso recebeu o conceito 1.
Quanto à Modalidade de Ensino, 49 dos 68 cursos na região Sudeste ( 72,1 \% ) foram oferecidos na Modalidade Presencial.
Para esses, a moda foi o conceito 4, com 24 cursos.
Os demais cursos Presenciais da região obtiveram os conceitos 1 ( um curso ), 2 ( sete cursos ), 3 ( dez cursos ) e 5 ( quatro cursos ).
Ficaram Sem Conceito ( SC ) três cursos.
Os 19 cursos \textbf{restantes} ( 27,9 \% ) da Modalidade Educação a Distância, na região Sudeste, tiveram distribuição \textbf{bimodal}, entre os conceitos 2 e 3, com sete cursos cada.
Os demais, nessa modalidade e região, obtiveram os conceitos 4 ( quatro cursos ) e 5 ( um curso ), não havendo cursos Sem Conceito ( SC ) ou com o conceito 1.
Dos 37 cursos participantes da região Sul, 14 foram de Instituições Privadas, correspondendo a 37,8 \% do total regional.
Desses, sete cursos obtiveram conceito 3, o conceito modal.
Nessa combinação de Categoria Administrativa e Grande Região, cinco cursos receberam o conceito 2, e um curso recebeu o conceito 4.
Um curso ficou Sem Conceito ( SC ), e nenhum recebeu os conceitos 1 ou 5.
Já as Instituições \textbf{Públicas} na região Sul concentraram 23 dos 37 cursos participantes ( 62,2 \% dos cursos da região ).
O conceito modal 3 foi atribuído a 12 cursos.
Os demais receberam os conceitos 2 ( dois cursos ), 4 ( seis cursos ) e 5 ( três cursos ).
Nenhum curso da região Sul oferecido por IES Públicas ficou Sem Conceito ( SC ) ou recebeu o conceito 1.
Na região Sul, 26 cursos ( 70,3 \% ) foram oferecidos na Modalidade Presencial, tendo sido o 3 ( 15 cursos ) o conceito modal nessa modalidade.
Aos demais foram atribuídos os conceitos 2 ( um curso ), 4 ( sete cursos ) e 5 ( três cursos ).
Nenhum curso ficou Sem Conceito ( SC ) ou recebeu o conceito 1.
Já dos 11 cursos ( 29,7 \% ) oferecidos nessa região na Modalidade a Distância, o conceito modal 2 foi obtido por seis cursos, e quatro cursos obtiveram o conceito 3.
Um curso da região Sul na Modalidade a Distância ficou Sem Conceito ( SC ), e nenhum recebeu os conceitos 1, 4 ou 5.
Na região Centro-Oeste, quatro cursos dos 31 participantes foram de Instituições Privadas ( 12,9 \% em termos regionais ).
Desses, dois cursos receberam o conceito modal 3.
Dos demais cursos, um recebeu o conceito 2, e um recebeu o conceito 4.
Nenhum curso, nessa Modalidade e região, ficou Sem Conceito ( SC ) ou recebeu os conceitos 1 e 5.
Quanto aos 27 cursos oferecidos por Instituições \textbf{Públicas} na região Centro Oeste ( 87,1 \% do total regional ), 13 cursos apresentaram conceito 2, o conceito modal.
Dos demais, dois receberam o conceito 1, oito receberam o conceito 3, três receberam o conceito 4, e um recebeu o conceito 5.
Nenhum curso ficou Sem Conceito ( SC ) nessa combinação de categoria e região.
Apenas três cursos na região Centro-Oeste foram oferecidos na Modalidade Educação a Distância, e o conceito modal foi o 3, obtido por dois cursos, tendo o outro curso obtido o conceito 2.
Nenhum curso ficou Sem Conceito ( SC ) nessa combinação de categoria e região, e nenhum recebeu os conceitos 1, 4 ou 5.
Em relação à Educação Presencial, foram 28 cursos no total, com o conceito modal 2 sendo obtido por 13 cursos ; o conceito 1 foi alocado a dois cursos, o 3 a oito cursos, o 4 a quatro cursos e o 5 a um curso.
No Gráfico 5.2a, apresenta -se a distribuição do Conceito Enade dos cursos de Licenciatura, segundo a Categoria Administrativa da IES.
Os cursos em IES Públicas ( linha azul ) apresentam uma \textbf{poligonal} mais à direita do que os em IES Privadas ( linha verde ) e, consequentemente, uma distribuição de Conceito Enade com valores maiores.
Na Tabela 5.2b, é apresentada a distribuição dos cursos dos cursos de Bacharelado de Geografia participantes do Enade/2021, por Categoria Administrativa e Modalidade de Ensino, de acordo com os conceitos por eles alcançados, segundo as Grandes Regiões brasileiras.
Dos 70 cursos participantes, 61 ( 87,1 \% ) eram \textbf{ministrados} em Instituições \textbf{Públicas}, e nove ( 12,9 \% ), em Instituições Privadas.
Quanto à Modalidade de Ensino, quatro ( 5,7 \% ) cursos eram a Distância, e 66 ( 94,3 \% ), Presenciais.
De acordo com as informações da Tabela 5.2b, em termos nacionais, entre Instituições \textbf{Públicas} e Privadas, observa -se que todos os cinco cursos avaliados com conceito 5 foram oferecidos em IES Públicas.
Dos 61 cursos participantes de IES Públicas, o conceito 3 foi o valor modal, atribuído a 24 cursos ( 39,3 \% ).
Dos demais cursos de Instituições \textbf{Públicas} participantes, três cursos ( 4,9 \% ) foram avaliados com o conceito 1 ; 11 cursos ( 18,0 \% ), com o conceito 2 ; 18 cursos ( 29,5 \% ), com o conceito 4 ; e cinco, ( 8,2 \% ), com o conceito 5.
Nenhum curso ficou Sem Conceito ( SC ) nessa categoria.
Na rede privada, os conceitos 3 e 4 foram atribuídos a três cursos cada ( \textbf{bimodal} ), correspondendo, juntos, a 66,6 \% do total.
Dos demais cursos participantes, dois ( 22,2 \% ) receberam o conceito 2, e um curso ( 11,1 \% ) ficou Sem Conceito ( SC ).
Dos 66 cursos ( 94,3 \% ) participantes da Modalidade Presencial, o conceito 3 foi o valor modal, atribuído a 25 cursos ( 37,9 \% ).
Dos demais cursos nessa Modalidade de Ensino, três cursos ( 4,5 \% ) receberam o conceito 1 ; 12 cursos ( 18,2 \% ) receberam o conceito 2 ; 21 cursos ( 31,8 \% ), o conceito 4 ; e cinco cursos ( 7,6 \% ), o conceito 5.
Nessa Modalidade de Ensino, nenhum curso ficou Sem Conceito ( SC ).
Dos quatro cursos a Distância informados, o conceito modal 3 foi atribuído a dois cursos, e o conceito 2, a um curso.
Houve ainda, um curso Sem Conceito ( SC ).
Na análise por região, observa -se que, na região Norte, todos os cursos foram mantidos por Instituições \textbf{Públicas} e eram Presenciais.
Três cursos receberam o conceito modal 3 ; o conceito 1 foi atribuído a um curso, e o conceito 2, a dois cursos.
Todos os cursos oferecidos na região Nordeste foram Presenciais, apenas um curso ( 6,3 \% do total da região ) dos 16 participantes era da rede privada, e esse obteve o conceito 2.
As Instituições \textbf{Públicas} da região Nordeste participaram com 15 cursos ( 93,8 \% do total da região ).
Desses, seis obtiveram conceito 2, o valor modal, e os demais cursos receberam os conceitos 3 ( quatro cursos ) e 4 ( cinco cursos ).
Nenhum curso ficou Sem Conceito ( SC ) ou recebeu os conceitos 1 ou 5.
Na região Sudeste, a proporção de cursos da rede privada foi de 21,7 \%, correspondendo a cinco dos 23 cursos participantes.
Nessa categoria e região, houve distribuição \textbf{bimodal} entre os conceitos 3 e 4 ( dois cursos cada ), e um curso obteve o conceito 2.
Para os 18 cursos oferecidos por Instituições \textbf{Públicas} na região Sudeste, a distribuição também foi \textbf{bimodal}, com os conceitos 3 e 4 sendo atribuídos a sete cursos cada.
Nessa região, um curso de IES Públicas recebeu o conceito 1, e três receberam o conceito 5.
Nenhum curso ficou Sem Conceito ( SC ) ou foi avaliado com o conceito 2.
Também na região Sudeste, a maioria dos cursos foi oferecida na Modalidade Presencial : 21 ( 91,3 \% ) dos 23 cursos.
Para esses, a moda foi o conceito 4, com nove cursos.
Os demais cursos presenciais da região receberam os conceitos 1 ( um curso ), 3 ( oito cursos ) e 5 ( três cursos ).
Na região Sudeste, os dois cursos ( 8,7 \% ) oferecidos na Modalidade Educação a Distância obtiveram os conceitos 2 e 3.
Na região Sul, três dos 16 cursos participantes, 18,8 \% do total regional, foram de Instituições Privadas.
Desses, um recebeu o conceito 3, outro recebeu o conceito 4 e outro ficou Sem Conceito ( SC ).
As Instituições \textbf{Públicas} na região Sul participaram com 13 cursos ( 81,3 \% dos cursos da região ), sendo que sete desses cursos receberam o conceito modal 3, cinco cursos receberam o conceito 4, e um recebeu o conceito 5.
Nenhum curso ficou Sem Conceito ( SC ) ou recebeu os demais conceitos.
A maioria dos cursos na região Sul foi oferecida na Modalidade Presencial : 14 ( 87,5 \% ) dos 16 cursos.
Para esses, a moda foi o conceito 3, com sete cursos.
Os demais cursos presenciais da região obtiveram os conceitos 4 ( seis cursos ) e 5 ( um curso ).
Na região Sul, dos dois cursos ( 12,5 \% ) oferecidos na Modalidade Educação a Distância, um obteve o conceito 3 e o outro ficou Sem Conceito ( SC ).
Na região Centro-Oeste, todos os nove cursos foram oferecidos por Instituições \textbf{Públicas} e eram Presenciais.
A distribuição dos conceitos, nessa região, foi \textbf{bimodal}, com três cursos recebendo o conceito 2, e três, o conceito 3.
Os demais conceitos nessa região foram 1, 4 e 5, atribuídos a um curso cada.
Nenhum curso ficou Sem Conceito ( SC ) nessa combinação de categoria e região.
No Gráfico 5.2b, apresenta -se a distribuição do Conceito Enade dos cursos de Bacharelado, segundo a Categoria Administrativa da IES.
Os cursos em IES Públicas ( linha azul ) apresentam uma \textbf{poligonal} mais à direita do que os em IES Privadas ( linha verde ) e, consequentemente, uma distribuição de Conceitos Enade com valores maiores.
Como nenhum curso de IES Privadas alcançou o conceito 5, a linha \textbf{poligonal} correspondente ( em verde ) alcança 100 \% no conceito 4. 5.3 CONCEITOS POR ORGANIZAÇÃO ACADÊMICA E GRANDE REGIÃO Na Tabela 5.3a, encontra -se a distribuição dos conceitos atribuídos aos cursos de Licenciatura participantes do Enade/2021, na Área de Geografia, por Organização Acadêmica, segundo as Grandes Regiões brasileiras.
Dos 229 cursos de Geografia participantes, 193 eram oferecidos em Universidades ; 13, em Centros Universitários ; nove, em Faculdades ; e 14, em CEFET/IF.
Essa distribuição corresponde a, respectivamente, 84,3 \%, 5,7 \%, 3,9 \% e 6,1 \% dos cursos.
De acordo com os dados apresentados, de todos os 11 cursos avaliados com conceito 5, dez eram vinculados a Universidades.
Os cursos vinculados a esse tipo de Organização Acadêmica tiveram o conceito 3 como conceito modal, atribuído a 71 cursos ( 36,8 \% ).
Os demais cursos vinculados a Universidades receberam os conceitos 1 ( oito cursos ), 2 ( 58 cursos ), 4 ( 43 cursos ) e conceito 5 ( dez cursos, como já mencionado ).
Três cursos ficaram Sem Conceito ( SC ).
Dos 13 cursos em Centros Universitários, o conceito modal foi 2, atribuído a seis cursos ( 46,2 \% ).
Dois cursos vinculados a esse tipo de Organização Acadêmica ficaram Sem Conceito ( SC ), e os outros cursos receberam os conceitos 3 ( quatro cursos ) e 5 ( um curso ).
Dos nove cursos mantidos por Faculdades, quatro ( 44,4 \% ) receberam conceito 2, o conceito modal.
O segundo conceito mais frequentemente atribuído aos cursos oferecidos por essa categoria de IES foi o conceito 3, recebido por três cursos.
Os dois cursos \textbf{restantes} receberam os conceitos 1 e 4.
Nenhum curso oferecido em Faculdades ficou Sem Conceito ( SC ) ou recebeu o conceito 5.
Já dos 14 cursos oferecidos por CEFET/IF, seis receberam conceito 4, valor modal, e os demais receberam os conceitos 2 e 3 ( quatro cursos cada ).
Considerando -se, separadamente, as regiões brasileiras, verifica -se que, na região Norte, as Universidades participaram com 17 dos 18 cursos avaliados, e o conceito modal foi o 2 ( oito cursos ).
Os demais receberam os conceitos 1 ( três cursos ), 3 ( quatro cursos ) e 4 ( um curso ).
Um curso ficou Sem Conceito ( SC ), e nenhum recebeu o conceito 5 nessa combinação de categoria e região.
Os CEFET/IF da região Norte foram representados por um curso, o qual obteve o conceito 2.
Na região Norte, não houve cursos oferecidos em Centros Universitários nem em Faculdades.
Na região Nordeste, as Universidades participaram com 63 dos 75 cursos da Área de Geografia ( Licenciatura ) oferecidos.
Nenhum desses 63 cursos ficou Sem Conceito ( SC ).
O conceito modal 3 foi atribuído a 30 cursos.
O segundo conceito mais frequente foi o 2 ( 23 cursos ).
Os dez cursos \textbf{restantes} obtiveram os conceitos 1 ( dois cursos ), 4 ( seis cursos ) e 5 ( dois cursos ).
Nenhum curso oferecido por Universidades no Nordeste ficou Sem Conceito ( SC ).
As Faculdades foram representadas por quatro cursos na região Nordeste, dois dos quais receberam conceito 2, o conceito modal.
Cada um dos conceitos 1 e 3 foi atribuído a um curso.
Na região Nordeste, nenhum curso, mantido por Faculdades, ficou Sem Conceito ( SC ) ou recebeu os demais conceitos.
Os CEFET/IF da região Nordeste foram representados por oito cursos, e o conceito modal foi o 4, atribuído a quatro cursos.
Três cursos obtiveram o conceito 2, e um curso obteve o conceito 3.
Na região Nordeste, nenhum curso, mantido por CEFET/IF, ficou Sem Conceito ( SC ) ou recebeu os demais conceitos.
Na região Nordeste, não houve cursos oferecidos em Centros Universitários.
Na região Sudeste, as Universidades concentraram 53 dos 68 cursos de Geografia ( Licenciatura ) da região.
Dos cursos oferecidos em Universidades, nessa região, o conceito modal foi o 4, atribuído a 25 cursos.
Os demais cursos receberam os conceitos 1 ( um curso ), 2 ( 11 cursos ), 3 ( dez cursos ) e 5 ( quatro cursos ).
Dois cursos ficaram Sem Conceito ( SC ).
Os Centros Universitários participaram com sete cursos na região Sudeste, dos quais três obtiveram o conceito modal 3.
O conceito 2 foi atribuído a dois cursos, e o conceito 5, a um curso.
Um curso ficou Sem Conceito ( SC ) nessa combinação de Organização Acadêmica e região, e nenhum recebeu os conceitos 1 ou 4.
Já as Faculdades foram representadas por três cursos na região Sudeste, que receberam os conceitos 2, 3 e 4.
Dos cinco cursos participantes vinculados a CEFET/IF da região Sudeste, três receberam o conceito modal 3, e dois, o conceito 4.
Dos 37 cursos da região Sul, 31 foram mantidos por Universidades, para os quais o conceito modal foi o 3, atribuído a 18 cursos.
Os demais cursos receberam os conceitos 2 ( três cursos ), 4 ( sete cursos ) e 5 ( três cursos ).
Nesse tipo de organização, nenhum curso da região Sul ficou Sem Conceito ( SC ) ou recebeu o conceito 1.
Dos cinco cursos participantes de Centros Universitários da região Sul, um ficou Sem Conceito ( SC ), três receberam o conceito 2 ( valor modal ), e um recebeu o conceito 3.
Já o único curso vinculado a Faculdades na região Sul recebeu o conceito 2.
Na região Sul, não houve cursos oferecidos em CEFET/IF.
Na região Centro-Oeste, 29 dos 31 cursos oferecidos foram mantidos por Universidades.
Para esse tipo de Organização Acadêmica, o conceito modal foi o 2, atribuído a 13 cursos.
Os demais cursos nessa combinação de Organização Acadêmica e região receberam os conceitos 1 ( dois cursos ), 3 ( nove cursos ), 4 ( quatro cursos ) e cinco ( um curso ), e nenhum ficou Sem Conceito ( SC ).
Os Centros Universitários da região Centro-Oeste \textbf{contaram} com apenas um curso, que obteve conceito 2.
Já o único curso oferecido por Faculdades na região Centro-Oeste recebeu o conceito 3.
Na região Centro-Oeste, não houve cursos oferecidos em CEFET/IF.
No Gráfico 5.3a, apresenta -se a distribuição do Conceito Enade dos cursos de Licenciatura, segundo a Organização Acadêmica da IES.
Os cursos em Universidades ( linha azul ) e CEFET/IF ( linha laranja ) apresentam uma \textbf{poligonal} mais à direita do que as demais, \textbf{denotando} uma melhor distribuição de Conceitos Enade.
No outro extremo, encontram -se os cursos mantidos por Faculdades ( linha vermelha ) e por Centros Universitários ( linha verde ), apresentando as respectivas \textbf{poligonais} mais à esquerda e, consequentemente, uma distribuição de Conceitos Enade com menores valores.
Como nenhum curso de Faculdades ( linha vermelha ) e de CEFET/IF ( linha laranja ) alcançou o conceito 5, as linhas \textbf{poligonais} correspondentes alcançaram 100 \% no conceito 4.
Na Tabela 5.3b, encontra -se a distribuição dos conceitos atribuídos aos cursos de Bacharelado participantes do Enade/2021, na Área de Geografia, por Organização Acadêmica, segundo as Grandes Regiões brasileiras.
Dos 70 cursos de Geografia participantes, 67 foram oferecidos em Universidades, e três, em Centros Universitários.
Essa distribuição corresponde a, respectivamente, 95,7 \% e 4,3 \% dos cursos.
De acordo com os dados apresentados, todos os cinco cursos avaliados com o conceito 5 eram vinculados a Universidades.
Os cursos vinculados a esse tipo de Organização Acadêmica tiveram o conceito 3 como o conceito modal, atribuído a 25 cursos ( 37,3 \% ).
Os demais cursos vinculados a Universidades receberam os conceitos 1 ( três cursos ), 2 ( 13 cursos ), 4 ( 21 cursos ) e conceito 5 ( cinco cursos, como já mencionado ), e nenhum curso ficou Sem Conceito ( SC ).
Dos cursos em Centros Universitários, o conceito modal foi o 3, atribuído a dois dos três cursos ( 66,7 \% ).
Apenas um curso ficou Sem Conceito ( SC ).
Considerando -se, separadamente, as regiões brasileiras, verifica -se que, na região Norte, todos os seis cursos avaliados foram vinculados a Universidades.
O conceito modal foi o 3, atribuído a três cursos.
Os demais obtiveram os conceitos 1 ( um curso ) e 2 ( dois cursos ).
Nenhum curso ficou Sem Conceito ( SC ) ou recebeu os demais conceitos.
Na região Nordeste, também todos os 16 cursos da Área de Geografia ( Bacharelado ) avaliados eram vinculados a Universidades.
O conceito modal foi o 2, com sete cursos, seguido do conceito 4, com cinco cursos, e do conceito 3, com quatro cursos.
Nenhum curso ficou Sem Conceito ( SC ) ou recebeu os demais conceitos.
Na região Sudeste, as Universidades concentraram 22 dos 23 cursos de Geografia ( Bacharelado ) da região.
Dos cursos oferecidos em Universidades, nessa região, o conceito modal foi o 4, atribuído a nove cursos.
Os demais cursos receberam os conceitos 1 ( um curso ), 2 ( um curso ), 3 ( oito cursos ) e 5 ( três cursos ).
Nenhum curso ficou Sem Conceito ( SC ).
Os Centros Universitários participaram com um curso na região Sudeste, que obteve o conceito 3.
Dos 16 cursos da região Sul, 14 foram mantidos por Universidades, para os quais o conceito modal foi 3, atribuído a sete cursos.
O segundo conceito mais frequente foi o 4, com seis cursos.
O curso \textbf{restante} recebeu o conceito 5.
Dos dois cursos participantes de Centros Universitários da região Sul, um ficou Sem Conceito ( SC ), e outro recebeu o conceito 3.
Na região Centro-Oeste, os nove cursos oferecidos foram mantidos por Universidades.
Para esse tipo de Organização Acadêmica, os conceitos 2 e 3 concentraram -se em 3 cursos cada ( distribuição \textbf{bimodal} ).
E os demais conceitos ( 1, 4 e 5 ) foram atribuídos a um curso cada, não havendo cursos Sem Conceito ( SC ) nessa combinação de Organização Acadêmica e região No Gráfico 5.3b, apresenta -se a distribuição do Conceito Enade dos cursos de Bacharelado, segundo a Organização Acadêmica da IES.
Nos Centros Universitários ( linha verde ), todos os cursos têm o conceito 3, já que a \textbf{poligonal} \textbf{cumulativa} alcançou 100 \% no conceito 3.
Os cursos em Universidades ( linha azul ) apresentaram uma \textbf{poligonal} mais à direita, \textbf{denotando} uma melhor distribuição de Conceitos Enade.

% matched lemmas: abscissa, bimodal, contar, cumulativo, denotar, equivaler, indicado, intermédio, ministrar, ordenado, panorama, poligonal, pontilhado, pública, restante, rosa, tabular, variável, visualizar
\end{textsample}
